\scanpage{82}\mspage{40}
peut prendre $=\lambda_0(\lambda)$ avec $\lambda\in\Lambda$ quelconque (pas néc$^{\text{t}}$ $\lambda\in\Lambda^*$). Mais sauf erreur,
ces trois types de matrices engendrent
$\mathrm{Sl}(2,\Lambda)$. On trouve ainsi (sauf erreur!)
% slanted
\margin{à vérifier}
que $u$ est l'identité sur $\widetilde S_M$ \deleted{tout} entier,
\deleted{vu que $\mathrm{Sl}(2,\Lambda)$ est transitif}
Comme $\mathrm{Gl}_\Lambda(M)$ est transitif sur l'ens. des
repères, on trouve ainsi :

Si $\Lambda^*$ engendre $\Lambda$ additivement, \deleted{alors} p.~ex. $\Lambda$ anneau
$\frac12$ local, alors tout autom. de la structure
\deleted{$\unsure{\widetilde P_M|\widetilde F_M}$} $(\widetilde S_M,\underset{\unsure{\mathfrak P(\mathfrak P_3(S))|\mathbb P(\mathbb P_3(S))|\mathcal P(\mathcal P_3(S))}}{\underset{\rotatebox[origin=c]{-90}{\(\in\)}}{\widetilde F_M}},T_\bullet:\Lambda^*\longrightarrow\circlearrowright_{\widetilde S_M})$ provient d'un
autom. unique de $M$.

Notons que pour une permutation
$u$ de $\widetilde S_M=M^*$, dire qu'elle commute
à l'antipodisme et commute à la structure
\unsure{triangulée|triangulation|triangulaire} signifie que l'on a
\[
\left\{\begin{aligned}
\text{a)}\ u(-x)&=-u(x)&&x\in M^*\\
\text{b)}\ u(x+y)&=u(x)+u(y)&&
\begin{aligned}&x,y\,\deleted{\unsure{x+y}}\in M^*\\[-2pt]&x\wedge y\in(\det M)^*\end{aligned}
\end{aligned}\right. \tag*{(21)}
\]
et dire qu'elle commute de plus à
l'action de $\Lambda^*$ signifie qu'on a, au
lieu de $(21\text{ a})$, la condition plus forte
\[
u(\lambda x)=\lambda u(x)\qquad x\in M^*,\ \lambda\in\Lambda^*. \tag*{(22)}
\]
% archive sheet number: 81
